\section{Parameter Tuning by Taguchi Orthogonal Array}
\label{sec:taguchi}

We replace empirical parameter choices by a Taguchi L9$(3^4)$ orthogonal-array design~\cite{montgomery,alrafei} over four factors at three levels---population $N\in\{30,50,70\}$, iteration budget $T\in\{100,300,500\}$, archive size $\in\{50,100,150\}$, and L\'evy exponent $\beta\in\{1.3,1.5,1.7\}$---with the larger-the-better signal-to-noise (S/N) ratio of the hypervolume as response, each of the nine configurations replicated over 12 seeds disjoint from the evaluation seeds.

Averaging the S/N ratio per factor level (Table~\ref{tab:response}, Figure~\ref{fig:taguchi}) gives an unambiguous ranking: \emph{population} ($\Delta=0.318$~dB) and \emph{iteration budget} ($\Delta=0.155$~dB) dominate, while $\beta$ ($0.108$) and archive size ($0.046$) lie within run-to-run noise; the population effect grows toward the top of its range while the iteration effect saturates past $T{=}300$. The additive model predicts the optimum $N{=}70$, $T{=}500$ at S/N $\approx109.77$~dB; a confirmation run gives $109.66$~dB ($\HV=304{,}391$), validating the model. The 30-run study adopts $N{=}50$, $T{=}500$, archive $100$, $\beta{=}1.5$, whose confirmation gives S/N $109.56$~dB---within $1.2\,\%$ hypervolume of the L9 optimum (its sole difference is $N{=}50$ vs.\ $70$), a deliberately conservative, compute-economical operating point ($T{=}500$ far exceeds the iteration-135 convergence point).

\begin{table}[t]
\centering
\caption{Response table: mean S/N (dB) per factor level, range $\Delta$, and rank. Larger $\Delta$ means a more influential factor; the highest-S/N level in each row is in bold. Each $\Delta$ is computed from unrounded S/N values.}
\label{tab:response}
\small
\begin{tabular}{lccccc}
\toprule
\textbf{Factor} & \textbf{Level 1} & \textbf{Level 2} & \textbf{Level 3} & $\boldsymbol{\Delta}$ & \textbf{Rank}\\
\midrule
A: population $N$        & 109.314 & 109.443 & \textbf{109.632} & 0.318 & 1\\
B: iteration budget $T$  & 109.372 & 109.489 & \textbf{109.528} & 0.155 & 2\\
C: archive size          & 109.461 & \textbf{109.487} & 109.441 & 0.046 & 4\\
D: Lévy exponent $\beta$ & 109.462 & \textbf{109.518} & 109.410 & 0.108 & 3\\
\bottomrule
\end{tabular}
\end{table}

\begin{figure}[t]
\centering
\includegraphics[width=\textwidth]{taguchi_main_effects.pdf}
\caption{Main-effects plot of the S/N ratio. The dashed line is the grand mean (109.46~dB) and the red star the best level per factor; population and iteration budget show clear trends (population the steeper), while archive size and $\beta$ are essentially flat.}
\label{fig:taguchi}
\end{figure}
